% Подключать этот файл в main только после локального запуска RUN 04
% и замены placeholder-таблицы на реальные значения из artifacts/run_04/metrics.csv.

\subsection{Дополнительный эксперимент: ablation признаков и снижение размерности}

Для углубления анализа был проведён дополнительный эксперимент RUN 04. Его цель состоит в том, чтобы проверить, зависит ли вывод о слабой добавочной силе FinBERT-признаков от высокой размерности embedding-пространства и от ограниченного набора рыночных признаков. В RUN 03 объединённый вектор признаков имел размерность

\[
x_t = [p_t;e_t] \in \mathbb{R}^{779},
\]

где $p_t \in \mathbb{R}^{11}$ --- рыночные признаки, а $e_t \in \mathbb{R}^{768}$ --- FinBERT embedding. При ограниченном числе наблюдений такая размерность может приводить к переобучению и снижать устойчивость результата. Поэтому в RUN 04 текстовые признаки дополнительно проверяются после снижения размерности методом главных компонент.

Для FinBERT embeddings строятся сжатые представления:

\[
e_t^{(k)} = PCA_k(e_t), \qquad e_t^{(k)} \in \mathbb{R}^{k}, \qquad k \in \{16,32,64\}.
\]

PCA обучается только на train subset, после чего то же преобразование применяется к validation и test. Это необходимо для предотвращения утечки информации из будущих периодов в preprocessing.

В RUN 04 сравниваются следующие группы признаков:

\begin{table}[htbp]
\centering
\begin{tabular}{ll}
\toprule
Группа & Состав \\
\midrule
\texttt{majority} & наиболее частый класс в train \\
\texttt{previous\_direction} & прогноз $\hat y_t = y_{t-1}$ \\
\texttt{prices\_raw} & исходные price/inc признаки \\
\texttt{returns\_rolling} & доходности, rolling mean, rolling volatility, residuals \\
\texttt{volume\_rolling} & log-volume, relative volume, volume z-score \\
\texttt{market\_full} & расширенные рыночные признаки \\
\texttt{finbert\_full} & полный FinBERT embedding \\
\texttt{finbert\_pca\_k} & PCA-сжатые FinBERT embeddings \\
\texttt{market\_full\_finbert\_pca\_k} & market features + PCA-сжатые embeddings \\
\bottomrule
\end{tabular}
\caption{Группы признаков в RUN 04}
\end{table}

Добавочная сила текстовых признаков в этом эксперименте определяется как

\[
\Delta_{text}^{M}(k) = M(X_{market}, E_k) - M(X_{market}),
\]

где $M$ --- выбранная метрика качества, $X_{market}$ --- расширенная группа рыночных признаков, а $E_k$ --- PCA-сжатые FinBERT-признаки размерности $k$.

% После запуска RUN 04 заменить таблицу ниже на реальные результаты из artifacts/run_04/metrics.csv.

\begin{table}[htbp]
\centering
\begin{tabular}{llrrrr}
\toprule
Model & Feature group & Accuracy & Balanced acc. & F1 & ROC-AUC \\
\midrule
PLACEHOLDER & majority & -- & -- & -- & -- \\
PLACEHOLDER & market\_full & -- & -- & -- & -- \\
PLACEHOLDER & finbert\_pca\_32 & -- & -- & -- & -- \\
PLACEHOLDER & market\_full\_finbert\_pca\_32 & -- & -- & -- & -- \\
\bottomrule
\end{tabular}
\caption{Результаты RUN 04 после локального запуска}
\end{table}

После получения результатов RUN 04 необходимо сравнить \texttt{market\_full} с \texttt{market\_full\_finbert\_pca\_k}. Если $\Delta_{text}^{M}(k)>0$ для нескольких метрик, то PCA-сжатые FinBERT-признаки дают добавочную предсказательную силу. Если же улучшение не наблюдается, то вывод RUN 03 усиливается: даже после снижения размерности и расширения рыночного baseline устойчивая добавочная сила текстовых признаков не подтверждается.
